% !TeX root = ../main_ustc.tex
\chapter{面向柔性接触任务的折扣帕累托力--位协同控制方法}
\label{chap:flexible_pareto_contact}

第三章在第\ref{chap:preliminaries}章预备知识的基础上，研究了刚性几何约束操作中的人机意图仲裁问题。该类任务的核心在于判断操作者监督意图，并在机器人自主目标与人侧监督目标之间分配控制权。然而，当机器人末端与柔性环境持续接触时，控制问题发生了本质变化。即使操作者意图保持不变，柔性环境刚度、阻尼、接触深度和扫描方向的变化也会导致位置误差与接触力误差相互耦合。此时，控制器不仅要跟踪期望轨迹，还要保证接触力位于安全范围内，并避免由 远心运动几何约束引起的姿态变化放大接触风险。

本章面向柔性接触恒力扫描任务，提出一种交互力安全裕度驱动的折扣帕累托力--位协同控制方法。为避免控制器仅停留在工程实现层面，本章首先从机械臂关节空间动力学出发，推导任务空间动力学模型，并在此基础上建立含残差项的力--位增广系统。与第三章的人机仲裁参数 $\alpha_{\mathrm{HR}}$ 不同，本章使用 $\alpha_{\mathrm{FP}}$ 表示力--位优先级。其值越大，控制器越偏向位置跟踪；其值越小，控制器越偏向接触力调节。为了提高理论表达的完整性，本章首先建立含残差项的力--位增广非线性系统，然后在折扣型多目标代价函数下求解帕累托反馈控制律，并进一步分析残差扰动、刚度失配和动态优先级变化对性能与稳定性的影响。本章沿用第\ref{chap:preliminaries}章中的接触建模和鲁棒性分析记号，使用 $\xi$ 表示力--位增广状态，使用 $\alpha_{\mathrm{FP}}$ 表示位置目标与力目标之间的优先级。

\section{问题描述与系统建模}
\label{sec:ch4_problem}

本节先给出柔性接触扫描任务的目标，再由机械臂动力学过渡到局部力--位增广模型，并补充环境参数估计方法，为后续控制器求解提供统一状态与参数基础。

\paragraph{任务场景与控制目标。}
\label{subsec:ch4_task}

考虑从端机器人携带细长工具通过固定几何约束点与柔性环境接触，并沿给定路径完成扫描。设 $x(t)\in\R^n$ 为工具接触点位置，$x_d(t)\in\R^n$ 为期望扫描轨迹，$f_e(t)\in\R^n$ 为环境作用于工具的接触力，$f_d(t)\in\R^n$ 为期望接触力。柔性环境可表示为组织表面、软质工件、弹性检测对象或可变刚度仿真接触面。扫描过程中，工具需要保持适当压入。接触力过大可能造成环境损伤或机器人冲击，接触力过小则可能导致脱离接触，使扫描失效。

设受控的标量接触力为
\begin{equation}
  y_f(t)=S_f f_e(t),
  \label{eq:ch4_yf}
\end{equation}
其中 $S_f\in\R^{1\times n}$ 为力方向选择矩阵。给定安全力范围 $[F_{\min},F_{\max}]$ 和期望力 $F_d$，满足
\begin{equation}
  F_{\min}<F_d<F_{\max} .
  \label{eq:ch4_force_band}
\end{equation}
本章希望控制 $y_f(t)$ 尽可能接近 $F_d$，同时避免越过安全上下界。与第三章不同，这里的关键权衡不再是“人侧目标与机器人侧目标”，而是“位置目标与力安全目标”。

\begin{problem}[柔性接触下的位置--力协同控制问题]
\label{prob:ch4_force_position}
考虑机器人末端工具在几何约束下与柔性环境保持接触并沿期望轨迹扫描的任务。设计反馈控制律 $u(t)$ 与力--位优先级参数 $\alpha_{\mathrm{FP}}(t)$，使系统在存在接触模型误差和残差扰动时同时满足以下目标：
\begin{enumerate}[label=(\arabic*)]
  \item 位置误差 $e_p(t)=x(t)-x_d(t)$ 保持有界并尽可能减小；
  \item 力误差 $e_f(t)=f_e(t)-f_d(t)$ 保持有界并尽可能减小；
  \item 受控接触力 $y_f(t)$ 尽可能位于 $[F_{\min},F_{\max}]$ 内；
  \item 几何约束误差 $e_{\mathrm{gc}}(t)$ 保持在允许范围内；
  \item 在线控制计算满足实时执行要求。
\end{enumerate}
\end{problem}

\begin{remark}
本章中的 $\alpha_{\mathrm{FP}}$ 与第三章中的 $\alpha_{\mathrm{HR}}$ 具有不同含义。$\alpha_{\mathrm{HR}}$ 表示人机目标权重，$\alpha_{\mathrm{FP}}$ 表示位置目标与力目标之间的优先级。两者不应在同一公式中混用。
\end{remark}


\paragraph{任务空间动力学模型。}
\label{subsec:ch4_robot_dynamics}

柔性接触控制的设计对象是从端机械臂在任务空间内的接触点运动。参考安全人机协作最优运动规划中的建模方式，先从机械臂关节空间动力学出发，再将其转换到任务空间\cite{li2025optimal}。设机械臂具有 $n_q$ 个关节，其动力学为
\begin{equation}
  M_q(q)\ddot q+C_q(q,\dot q)\dot q+G_q(q)=\tau,
  \label{eq:ch4_joint_dynamics}
\end{equation}
其中 $q,\dot q,\ddot q\in\R^{n_q}$ 分别为关节位置、速度和加速度，$M_q(q)$ 为惯性矩阵，$C_q(q,\dot q)$ 为科氏力与离心力矩阵，$G_q(q)$ 为重力项，$\tau$ 为关节力矩输入。

设工具接触点或等效控制点位置为 $p\in\R^n$，正运动学为
\begin{equation}
  p=\Phi(q),
  \label{eq:ch4_forward_kinematics}
\end{equation}
其速度关系为
\begin{equation}
  \dot p=J(q)\dot q .
  \label{eq:ch4_task_velocity}
\end{equation}
在 $J(q)$ 行满秩的工作区间内，有
\begin{equation}
  \dot q=J^+(q)\dot p,
  \quad
  \ddot q=\dot J^+(q)\dot p+J^+(q)\ddot p .
  \label{eq:ch4_qdot_qddot}
\end{equation}
将\eqref{eq:ch4_qdot_qddot}代入\eqref{eq:ch4_joint_dynamics}并左乘 $J^{+\T}(q)$，得到任务空间动力学
\begin{equation}
  M_p(q)\ddot p+C_p(q,\dot q)\dot p+G_p(q)=\mu,
  \label{eq:ch4_task_dynamics}
\end{equation}
其中
\begin{align}
  M_p(q)&=J^{+\T}(q)M_q(q)J^+(q),
  \label{eq:ch4_Mp}\\
  C_p(q,\dot q)&=J^{+\T}(q)\bigl(M_q(q)\dot J^+(q)+C_q(q,\dot q)J^+(q)\bigr),
  \label{eq:ch4_Cp}\\
  G_p(q)&=J^{+\T}(q)G_q(q),
  \label{eq:ch4_Gp}\\
  \mu&=J^{+\T}(q)\tau.
  \label{eq:ch4_mu}
\end{align}
若局部常秩条件成立，上式可进一步转换为含 $\dot J(q)$ 的常见表达。本文采用\eqref{eq:ch4_Cp}是为了避免在一般伪逆情形下省略 $\dot J^+(q)$ 所引入的附加项。

定义机器人任务空间状态
\begin{equation}
  \zeta(t)=\begin{bmatrix}p^{\T}(t)&\dot p^{\T}(t)\end{bmatrix}^{\T} .
  \label{eq:ch4_zeta_state}
\end{equation}
由\eqref{eq:ch4_task_dynamics}可得任务空间非线性状态模型
\begin{equation}
  \dot\zeta=F_m(\zeta)\zeta+B_m(\zeta)u_m,
  \label{eq:ch4_nonlinear_task_model}
\end{equation}
其中
\begin{equation}
  F_m(\zeta)=
  \begin{bmatrix}
    0&I\\
    0&-M_p^{-1}(q)C_p(q,\dot q)
  \end{bmatrix},
  \quad
  B_m(\zeta)=
  \begin{bmatrix}
    0\\
    M_p^{-1}(q)
  \end{bmatrix},
  \quad
  u_m=\mu-G_p(q).
  \label{eq:ch4_Fm_Bm}
\end{equation}
该模型体现了机械臂惯性、速度相关项和重力项对任务空间接触运动的影响。柔性接触控制器并不直接对\eqref{eq:ch4_nonlinear_task_model}求解完整非线性哈密顿--雅可比--贝尔曼方程，而是利用内环补偿和局部线性化得到可实时求解的增广误差模型。由任务空间非线性模型到局部增广模型之间的未建模项，将在后续统一并入残差 $d(\xi,t)$。

\subsection{力--位增广系统与环境估计}
\label{subsec:ch4_augmented_system}

在任务空间动力学\eqref{eq:ch4_task_dynamics}及其状态模型\eqref{eq:ch4_nonlinear_task_model}基础上，柔性接触任务采用修正阻抗动力学作为控制器设计的局部等效模型。阻抗控制为机器人与环境交互提供了可解释的柔顺动力学结构\cite{hogan1985impedance}：
\begin{equation}
  M\ddot x+C\dot x+K_v(x-x_d)=\mu+f_e,
  \label{eq:ch4_impedance}
\end{equation}
其中 $M\succ0$ 为期望惯性矩阵，$C\succ0$ 为期望阻尼矩阵，$K_v\succ0$ 为基线虚拟刚度矩阵，$\mu$ 为机器人任务空间输入。基线虚拟刚度提供名义位置恢复作用，并改善后续增广误差系统的正则性。

柔性环境在局部接触区间内采用开尔文--沃伊特模型近似：
\begin{equation}
  f_e=-K_e(x-x_s)-B_e(\dot x-\dot x_s)+\Delta_f(x,\dot x,t),
  \label{eq:ch4_kv}
\end{equation}
其中 $x_s(t)$ 为局部接触面参考位置，$K_e\succeq0$ 和 $B_e\succeq0$ 分别为局部等效刚度与阻尼，$\Delta_f$ 表示接触模型误差。对于分段变刚度环境，$K_e$ 与 $B_e$ 可在不同扫描区域取不同常值；对于真实柔性环境，它们只是局部近似参数，并不要求精确已知。

\begin{assumption}[柔性接触模型局部有效性]
\label{ass:ch4_contact_model}
在每一个短时间接触区间内，柔性环境可由\eqref{eq:ch4_kv}近似描述。矩阵 $K_e$、$B_e$ 以及接触面参考 $x_s(t)$、$\dot x_s(t)$ 在所考虑的紧集内有界。模型误差 $\Delta_f(x,\dot x,t)$ 及其对时间的局部导数有界。
\end{assumption}

\paragraph{环境参数在线估计。}
\label{subsec:ch4_env_estimator}

上述控制器设计依赖局部环境参数。对于真实柔性环境，$K_e$ 和 $B_e$ 通常无法事先精确获得，并且可能随扫描位置、压入深度和材料非均匀性变化。因此，本章采用面向控制增益调度的在线环境估计方法，实时给出局部表观刚度估计 $\hat K_e$ 和阻尼估计 $\hat B_e$。该估计量不用于替代完整接触模型，而是作为离线增益表中的调度变量，使控制器能够根据环境软硬变化调整反馈增益。

沿受控接触方向定义单向压入量
\begin{equation}
  \delta(t)=\max\{0,x_s(t)-x_n(t)\},
  \label{eq:ch4_delta_contact}
\end{equation}
其中 $x_n(t)$ 为工具尖端在接触法向上的坐标。未接触时 $\delta=0$；发生压入时 $\delta>0$。在局部短时间区间内，受控力可近似写为标量开尔文--沃伊特形式
\begin{equation}
  y_f(t)=K_e\delta(t)+B_e\dot\delta(t)+\nu_f(t),
  \label{eq:ch4_scalar_kv}
\end{equation}
其中 $\nu_f(t)$ 表示测量噪声和模型误差。若直接对 $K_e$ 和 $B_e$ 同时做无约束最小二乘估计，速度项 $\dot\delta$ 的噪声容易被误解释为刚度突变，尤其是在恒力扫描中压入量通常只有毫米量级。为提高估计对控制的可用性，本章将刚度估计和阻尼估计分开处理：$\hat K_e$ 采用准静态表观刚度的低通估计，$\hat B_e$ 采用有界递推最小二乘作为辅助估计。

当接触力和压入量足够大时，定义表观刚度观测值
\begin{equation}
  K_{\mathrm{obs}}(t)=
  \sat\left(\frac{|y_f(t)|}{|\delta(t)|},K_{\min},K_{\max}\right),
  \quad |y_f|\ge F_{\mathrm{th}},\ |\delta|\ge\delta_{\min}.
  \label{eq:ch4_K_obs}
\end{equation}
若不满足接触判据，则保持上一时刻估计值不变，避免未接触或微小压入时的信息不足污染估计。刚度估计通过指数低通更新：
\begin{equation}
  \hat K_{e,k}=
  \sat\left(\hat K_{e,k-1}+\alpha_K\left(K_{\mathrm{obs},k}-\hat K_{e,k-1}\right),
  K_{\min},K_{\max}\right),
  \quad 0<\alpha_K<1 .
  \label{eq:ch4_K_ema}
\end{equation}
该式使 $\hat K_e$ 能够跟随分段刚度变化，同时抑制接触力瞬态和速度估计噪声造成的尖峰。

阻尼估计采用带遗忘因子的有界递推最小二乘。令
\begin{equation}
  \phi_k=\begin{bmatrix}\delta_k&\bar{\dot\delta}_k\end{bmatrix}^{\T},
  \quad
  \theta_k=\begin{bmatrix}K_{e,k}&B_{e,k}\end{bmatrix}^{\T},
  \label{eq:ch4_rls_phi_theta}
\end{equation}
其中 $\bar{\dot\delta}_k$ 为限幅后的压入速度。递推增益、参数和协方差更新为
\begin{align}
  L_k&=\frac{P_{k-1}\phi_k}{\lambda_R+\phi_k^{\T}P_{k-1}\phi_k},
  \label{eq:ch4_rls_gain}\\
  \theta_k&=\theta_{k-1}+L_k\left(y_{f,k}-\phi_k^{\T}\theta_{k-1}\right),
  \label{eq:ch4_rls_theta}\\
  P_k&=\frac{1}{\lambda_R}\left(P_{k-1}-L_k\phi_k^{\T}P_{k-1}\right),
  \label{eq:ch4_rls_P}
\end{align}
其中 $\lambda_R\in(0,1]$ 为遗忘因子。更新后对 $K_{e,k}$ 和 $B_{e,k}$ 施加物理上下界约束。实际反馈增益调度使用\eqref{eq:ch4_K_ema}中的 $\hat K_e$，而 $\hat B_e$ 主要作为接触模型诊断和阻尼辅助信息。这样既保留了开尔文--沃伊特模型的可解释性，又避免速度噪声主导刚度调度。

定义位置误差、速度误差和力误差为
\begin{equation}
  e_p=x-x_d,
  \quad e_v=\dot x-\dot x_d,
  \quad e_f=f_e-f_d .
  \label{eq:ch4_errors}
\end{equation}
为改善稳态力调节并避免积分饱和，引入力误差泄漏积分状态
\begin{equation}
  \dot\sigma_f=e_f-\lambda_f\sigma_f,
  \quad \lambda_f>0 .
  \label{eq:ch4_sigma}
\end{equation}
增广误差状态为
\begin{equation}
  \xi=\begin{bmatrix}e_p^{\T}&e_v^{\T}&e_f^{\T}&\sigma_f^{\T}\end{bmatrix}^{\T}\in\R^{4n} .
  \label{eq:ch4_xi}
\end{equation}

\begin{assumption}[残差项有界性]
\label{ass:ch4_residual}
将任务空间动力学线性化误差、内环补偿误差、参考轨迹变化、接触面运动、接触模型误差、几何映射误差、摩擦、离散化误差和测量噪声统一记为残差项 $d(\xi,t)$。存在非负常数 $\bar d_0$ 与 $\bar d_1$，使得
\begin{equation}
  \norm{d(\xi,t)}\le \bar d_0+\bar d_1\norm{\xi},\quad \forall t\ge0 .
  \label{eq:ch4_residual_bound_affine}
\end{equation}
若只在给定紧集 $\Omega_\xi$ 内分析局部性能，则可取常数 $\bar d>0$，使得
\begin{equation}
  \norm{d(\xi,t)}\le \bar d,
  \quad \xi\in\Omega_\xi .
  \label{eq:ch4_residual_bound_const}
\end{equation}
\end{assumption}

\begin{lemma}[含残差项的增广位置--力误差系统]
\label{lem:ch4_augmented_system}
在\cref{ass:ch4_contact_model}下，修正阻抗系统\eqref{eq:ch4_impedance}与局部开尔文--沃伊特模型\eqref{eq:ch4_kv}可写成增广误差系统
\begin{equation}
  \dot\xi=A_c\xi+B_c\mu+d(\xi,t),
  \label{eq:ch4_aug_system_mu}
\end{equation}
其中 $A_c$ 和 $B_c$ 由 $M,C,K_v,K_e,B_e,\lambda_f$ 的局部标称值确定。若采用前馈补偿 $\mu=u_{\mathrm{ff}}+u$，并将未补偿项并入残差，则控制器设计所使用的标称反馈模型为
\begin{equation}
  \dot\xi=A_c\xi+B_cu .
  \label{eq:ch4_nominal_model}
\end{equation}
\end{lemma}

\begin{proof}
由 $e_p=x-x_d$ 可得
\begin{equation}
  \dot e_p=\dot x-
  \dot x_d=e_v .
  \label{eq:ch4_ep_dot}
\end{equation}
由 $e_v=\dot x-\dot x_d$ 可得
\begin{equation}
  \dot e_v=\ddot x-\ddot x_d .
  \label{eq:ch4_ev_dot_def}
\end{equation}
根据修正阻抗动力学\eqref{eq:ch4_impedance}，有
\begin{equation}
  M\ddot x=-C\dot x-K_v(x-x_d)+\mu+f_e .
\end{equation}
两边左乘 $M^{-1}$，并代入 $\dot x=e_v+\dot x_d$、$x-x_d=e_p$、$f_e=e_f+f_d$，得到
\begin{equation}
  \dot e_v=-M^{-1}K_ve_p-M^{-1}Ce_v+M^{-1}e_f+M^{-1}\mu+\Delta_v(t),
  \label{eq:ch4_ev_dot}
\end{equation}
其中
\begin{equation}
  \Delta_v(t)=-M^{-1}C\dot x_d+M^{-1}f_d-\ddot x_d
  \label{eq:ch4_delta_v}
\end{equation}
为参考轨迹和期望力引起的仿射项。若该项可测或可计算，可由 $u_{\mathrm{ff}}$ 部分补偿；未补偿部分进入残差项。

接下来考虑力误差动态。由 $e_f=f_e-f_d$ 得
\begin{equation}
  \dot e_f=\dot f_e-\dot f_d .
  \label{eq:ch4_ef_dot_def}
\end{equation}
对\eqref{eq:ch4_kv}求导，在局部短时间区间内将 $K_e$ 和 $B_e$ 的变化归入残差，得到
\begin{equation}
  \dot f_e=-K_e(\dot x-\dot x_s)-B_e(\ddot x-\ddot x_s)+\Delta_{\dot f}(x,\dot x,t),
  \label{eq:ch4_fe_dot}
\end{equation}
其中 $\Delta_{\dot f}$ 包含接触参数变化和模型误差导数项。将 $\dot x=e_v+\dot x_d$ 以及由\eqref{eq:ch4_impedance}得到的 $\ddot x$ 代入\eqref{eq:ch4_fe_dot}，并整理关于 $e_p,e_v,e_f,\mu$ 的项，可得
\begin{equation}
  \dot e_f=B_eM^{-1}K_ve_p+(-K_e+B_eM^{-1}C)e_v-B_eM^{-1}e_f-B_eM^{-1}\mu+\Delta_f^{\mathrm{res}}(\xi,t),
  \label{eq:ch4_ef_dot}
\end{equation}
其中 $\Delta_f^{\mathrm{res}}$ 汇集 $\dot x_d,\ddot x_d,\dot f_d,x_s,\dot x_s,\ddot x_s$ 以及模型失配项。由\eqref{eq:ch4_sigma}有
\begin{equation}
  \dot\sigma_f=e_f-\lambda_f\sigma_f .
  \label{eq:ch4_sigma_dot_final}
\end{equation}
将\eqref{eq:ch4_ep_dot}、\eqref{eq:ch4_ev_dot}、\eqref{eq:ch4_ef_dot}与\eqref{eq:ch4_sigma_dot_final}依次堆叠，即可得到\eqref{eq:ch4_aug_system_mu}。若使用 $\mu=u_{\mathrm{ff}}+u$ 对可计算仿射项进行补偿，剩余反馈设计模型为\eqref{eq:ch4_nominal_model}，未补偿项并入 $d(\xi,t)$。引理得证。
\end{proof}

对于逐轴实现，可取单轴状态 $\xi_i=[e_{p,i},e_{v,i},e_{f,i},\sigma_{f,i}]^{\T}$，并写成
\begin{equation}
  \dot\xi_i=\bar A_i\xi_i+\bar b_i u_i+d_i(t),
  \label{eq:ch4_axis_state}
\end{equation}
其中
\begin{equation}
  \bar A_i=\begin{bmatrix}
  0&1&0&0\\
  -K_v/M&-C/M&1/M&0\\
  B_eK_v/M&-K_e+B_eC/M&-B_e/M&0\\
  0&0&1&-\lambda_f
  \end{bmatrix},
  \quad
  \bar b_i=\begin{bmatrix}0\\1/M\\-B_e/M\\0\end{bmatrix} .
  \label{eq:ch4_axis_matrices}
\end{equation}
该模型是后续离线增益表和在线反馈计算的实现基础。三维任务空间可由逐轴模型按块对角方式扩展。

\subsection{折扣多目标控制问题}
\label{subsec:ch4_discount_problem}

\begin{definition}[折扣型位置目标与力目标代价函数]
\label{def:ch4_discount_costs}
对位置跟踪目标和力调节目标，分别定义折扣型无限时域代价函数
\begin{equation}
  J_p(\xi_0,u)=\int_0^\infty \e^{-\gamma t}\left[\xi^{\T}Q_p\xi+u^{\T}R_pu\right]\dd t,
  \label{eq:ch4_Jp}
\end{equation}
\begin{equation}
  J_f(\xi_0,u)=\int_0^\infty \e^{-\gamma t}\left[\xi^{\T}Q_f\xi+u^{\T}R_fu\right]\dd t,
  \label{eq:ch4_Jf}
\end{equation}
其中 $\gamma>0$ 为折扣因子，$Q_p\succeq0$ 主要惩罚 $e_p$ 和 $e_v$，$Q_f\succeq0$ 主要惩罚 $e_f$ 和 $\sigma_f$，$R_p,R_f\succ0$ 为控制输入权重。
\end{definition}

折扣因子使近期误差和近期力风险具有更高权重，也使无限时域代价在存在扰动和局部模型近似时具有更好的数学性质。类似的折扣帕累托思路已被用于连续接触任务中的位置--力协调控制\cite{li2026pareto}。给定力--位优先级 $\alpha_{\mathrm{FP}}\in[\alpha_{\min},\alpha_{\max}]$，定义
\begin{equation}
  Q_\alpha=\alpha_{\mathrm{FP}}Q_p+\bigl(1-\alpha_{\mathrm{FP}}\bigr)Q_f,
  \label{eq:ch4_Qalpha}
\end{equation}
\begin{equation}
  R_\alpha=\alpha_{\mathrm{FP}}R_p+\bigl(1-\alpha_{\mathrm{FP}}\bigr)R_f .
  \label{eq:ch4_Ralpha}
\end{equation}
对应的标量化折扣代价为
\begin{equation}
  J_\alpha(\xi_0,u)=\int_0^\infty \e^{-\gamma t}\left[\xi^{\T}Q_\alpha\xi+u^{\T}R_\alpha u\right]\dd t .
  \label{eq:ch4_Jalpha}
\end{equation}

\begin{definition}[帕累托有效控制律]
\label{def:ch4_pareto}
若不存在另一个可容许控制律 $\tilde u$，使得 $J_p(\xi_0,\tilde u)\le J_p(\xi_0,u)$ 且 $J_f(\xi_0,\tilde u)\le J_f(\xi_0,u)$，并且至少一个不等式严格成立，则称控制律 $u$ 对双目标代价 $(J_p,J_f)$ 是帕累托有效的。
\end{definition}

在凸二次情形下，加权标量化代价\eqref{eq:ch4_Jalpha}的最优解是支持帕累托解。在线阶段通过调节 $\alpha_{\mathrm{FP}}(t)$，可以在位置跟踪和力调节之间移动工作点，而不必每个控制周期重新求解完整多目标优化问题。


\section{折扣帕累托控制器设计}
\label{sec:ch4_game_controller}

在前述增广模型和环境估计基础上，本节将位置跟踪与力调节写成两个参与方的折扣代价，并由帕累托权重给出统一反馈律。这样既保留位置--力权衡的可解释性，又便于在线根据接触安全状态选择控制增益。

\subsection{位置--力参与方与帕累托反馈律}
\label{subsec:ch4_players}

本章将位置跟踪和力调节视为两个合作参与方。位置参与方主要惩罚 $e_p,e_v$，使工具沿期望扫描路径运动；力参与方主要惩罚 $e_f,\sigma_f$，使接触力接近期望值并抑制稳态偏差。二者共享同一个控制输入 $u$，因此不能通过两个互不相关的控制回路分别求解。柔性环境刚度变化时，若单纯追求位置跟踪，可能导致力超调；若单纯追求力调节，则可能牺牲扫描轨迹覆盖质量。帕累托控制的作用正是显式描述这种任务级折中。

\begin{assumption}[标称系统可镇定与可检测]
\label{ass:ch4_stabilizable}
标称系统矩阵对 $(A_c,B_c)$ 可镇定。对任意 $\alpha_{\mathrm{FP}}\in[\alpha_{\min},\alpha_{\max}]\subset(0,1)$，矩阵对 $(Q_\alpha^{1/2},A_c)$ 可检测，并且 $R_\alpha\succ0$。
\end{assumption}

\begin{lemma}[标称系统下的折扣哈密顿--雅可比--贝尔曼方程]
\label{lem:ch4_hjb}
对于标称系统\eqref{eq:ch4_nominal_model}和折扣代价\eqref{eq:ch4_Jalpha}，若 $V_\alpha(\xi)$ 为最优值函数，则其满足折扣哈密顿--雅可比--贝尔曼方程
\begin{equation}
  0=\min_u\left[\xi^{\T}Q_\alpha\xi+u^{\T}R_\alpha u+\nabla V_\alpha^{\T}(A_c\xi+B_cu)-\gamma V_\alpha\right].
  \label{eq:ch4_hjb}
\end{equation}
\end{lemma}

\begin{proof}
由动态规划原理，从任意时刻 $t$ 开始的最优值函数等于当前瞬时代价与后续最优代价之和。对折扣型无限时域问题，在一个微小时间区间 $[t,t+\Delta t]$ 上展开，有
\begin{align}
  V_\alpha(\xi(t))=&\min_u\int_t^{t+\Delta t}\e^{-\gamma(\tau-t)}\ell_\alpha(\xi(\tau),u(\tau))\dd\tau \\
  &+\e^{-\gamma\Delta t}V_\alpha(\xi(t+\Delta t)),
\end{align}
其中 $\ell_\alpha(\xi,u)=\xi^{\T}Q_\alpha\xi+u^{\T}R_\alpha u$。对 $V_\alpha(\xi(t+\Delta t))$ 作一阶展开，并令 $\Delta t$ 趋近于零，整理可得\eqref{eq:ch4_hjb}。引理得证。
\end{proof}

\paragraph{帕累托反馈律求解。}
\label{subsec:ch4_pareto_law}

\begin{theorem}[标称系统下的折扣帕累托最优反馈律]
\label{thm:ch4_discount_lqr}
在\cref{ass:ch4_stabilizable}成立的条件下，给定固定 $\alpha_{\mathrm{FP}}\in[\alpha_{\min},\alpha_{\max}]$。若 $P_\alpha=P_\alpha^{\T}\succeq0$ 是折扣代数黎卡提方程
\begin{equation}
  A_c^{\T}P_\alpha+P_\alpha A_c-P_\alpha B_cR_\alpha^{-1}B_c^{\T}P_\alpha+Q_\alpha-\gamma P_\alpha=0
  \label{eq:ch4_discount_are}
\end{equation}
的稳定化解，则折扣标量化代价\eqref{eq:ch4_Jalpha}对应的最优反馈律为
\begin{equation}
  u_\alpha^*(\xi)=-R_\alpha^{-1}B_c^{\T}P_\alpha\xi .
  \label{eq:ch4_u_star}
\end{equation}
\end{theorem}

\begin{proof}
设二次型值函数为
\begin{equation}
  V_\alpha(\xi)=\xi^{\T}P_\alpha\xi .
\end{equation}
则 $\nabla V_\alpha=2P_\alpha\xi$。将其代入\eqref{eq:ch4_hjb}中的哈密顿函数，得到
\begin{equation}
  \mathcal{H}_\alpha(\xi,u)=\xi^{\T}Q_\alpha\xi+u^{\T}R_\alpha u+2\xi^{\T}P_\alpha A_c\xi+2\xi^{\T}P_\alpha B_cu-\gamma\xi^{\T}P_\alpha\xi .
\end{equation}
对 $u$ 求偏导，有
\begin{equation}
  \frac{\partial\mathcal{H}_\alpha}{\partial u}=2R_\alpha u+2B_c^{\T}P_\alpha\xi .
\end{equation}
由 $R_\alpha\succ0$ 可知哈密顿函数关于 $u$ 严格凸。令偏导为零，得到唯一极小点
\begin{equation}
  u_\alpha^*(\xi)=-R_\alpha^{-1}B_c^{\T}P_\alpha\xi .
\end{equation}
将该控制律代回\eqref{eq:ch4_hjb}，得到
\begin{equation}
  \xi^{\T}\left[A_c^{\T}P_\alpha+P_\alpha A_c-P_\alpha B_cR_\alpha^{-1}B_c^{\T}P_\alpha+Q_\alpha-\gamma P_\alpha\right]\xi=0 .
\end{equation}
由于该式对任意 $\xi$ 成立，括号中的矩阵必须为零，因此 $P_\alpha$ 满足\eqref{eq:ch4_discount_are}。定理得证。
\end{proof}

\paragraph{离线表与在线插值。}
\label{subsec:ch4_offline_online}

上述折扣帕累托控制律给出了固定 $\alpha_{\mathrm{FP}}$ 与局部接触参数下的最优反馈解。为了满足高频接触控制实时性，实际实现采用离线--在线结构。离线阶段在 $\alpha_{\mathrm{FP}}$ 与局部刚度估计 $\hat K_e$ 的二维网格上求解\eqref{eq:ch4_discount_are}，得到增益表
\begin{equation}
  K(\alpha_{\mathrm{FP}},\hat K_e)=R_\alpha^{-1}B_c^{\T}P_\alpha .
  \label{eq:ch4_gain_table}
\end{equation}
在线阶段只更新 $\alpha_{\mathrm{FP}}(t)$ 和环境估计器给出的 $\hat K_e(t)$，并通过插值得到当前反馈增益。这样，控制循环不需要在线求解哈密顿--雅可比--贝尔曼方程、黎卡提方程或优化问题。该实现顺序也明确区分了“控制器求解”“环境参数估计”和“优先级仲裁”：控制器给出每个候选优先级与局部刚度下的最优反馈律，环境估计器提供当前表观刚度，仲裁器只负责根据接触安全状态选择合适的优先级。

\section{安全仲裁与约束执行}
\label{sec:ch4_arbitration_execution}

折扣帕累托控制器给出了可调优先级下的位置--力反馈律，但优先级本身仍需依据接触风险在线确定。下面从力安全裕度出发构造仲裁参数，并说明在远心运动约束下如何将接触点参考转化为法兰空间执行命令。

\subsection{力安全裕度仲裁}
\label{subsec:ch4_force_margin}

在折扣帕累托控制器已经建立的基础上，需要进一步说明 $\alpha_{\mathrm{FP}}(t)$ 如何在线选择。旧的仲裁方式常将环境刚度作为模糊输入之一。然而，刚度是环境属性，不直接等价于安全风险。高刚度环境中只要压入很浅，接触力仍可能安全；低刚度环境中若压入过深，接触力也可能越界。因此，本章将交互力相对上下安全边界的裕度作为 $\alpha_{\mathrm{FP}}$ 调节的核心依据。

\begin{definition}[交互力安全裕度]
\label{def:ch4_force_margin}
给定受控力 $y_f(t)$ 和安全范围 $[F_{\min},F_{\max}]$，定义到下界和上界的距离为
\begin{equation}
  d_{\mathrm{lower}}(t)=y_f(t)-F_{\min},
  \quad d_{\mathrm{upper}}(t)=F_{\max}-y_f(t) .
  \label{eq:ch4_force_distances}
\end{equation}
归一化安全裕度定义为
\begin{equation}
  \rho_F(t)=\sat\left(\frac{\min\{d_{\mathrm{lower}}(t),d_{\mathrm{upper}}(t)\}}{(F_{\max}-F_{\min})/2},0,1\right) .
  \label{eq:ch4_rhoF}
\end{equation}
边界方向变量定义为
\begin{equation}
  s_F(t)=\sat\left(\frac{y_f(t)-F_d}{(F_{\max}-F_{\min})/2},-1,1\right) .
  \label{eq:ch4_sF}
\end{equation}
其中 $\rho_F(t)$ 越接近 $1$ 表示接触力越接近安全区间中部，越接近 $0$ 表示越接近上下边界；$s_F(t)>0$ 表示接触力偏大，$s_F(t)<0$ 表示接触力偏小。
\end{definition}

模糊输入包括力误差幅值、安全裕度和位置误差幅值。定义归一化力偏差
\begin{equation}
  r_e(t)=\frac{|y_f(t)-F_d|}{\varepsilon_f},
  \label{eq:ch4_re}
\end{equation}
其中 $\varepsilon_f>0$ 为允许力误差尺度。位置误差指标可取
\begin{equation}
  r_p(t)=\frac{\norm{e_p(t)}}{\varepsilon_p},
  \label{eq:ch4_rp}
\end{equation}
其中 $\varepsilon_p>0$ 为允许位置误差尺度。模糊推理输出基础优先级
\begin{equation}
  \alpha_0(t)=\mathcal{F}_{\mathrm{FP}}(r_e(t),\rho_F(t),r_p(t)).
  \label{eq:ch4_alpha0}
\end{equation}

力接近上界或下界时均需要增强力调节能力，因此安全修正应降低位置优先级。定义风险强度 $r_F(t)=1-\rho_F(t)$，并取
\begin{equation}
  \alpha_{\mathrm{safe}}(t)=\alpha_0(t)-k_{\mathrm{upper}}r_F(t)\max\{s_F(t),0\}-k_{\mathrm{lower}}r_F(t)\max\{-s_F(t),0\},
  \label{eq:ch4_alpha_safe}
\end{equation}
其中 $k_{\mathrm{upper}}>0$ 与 $k_{\mathrm{lower}}>0$ 为修正系数。当接触力接近上界时，控制器降低 $\alpha_{\mathrm{FP}}$ 以增强力误差调节；当接触力接近下界时，控制器同样降低 $\alpha_{\mathrm{FP}}$，以优先恢复目标接触力并避免脱离接触。最终优先级由限幅和一阶滤波得到：
\begin{equation}
  \alpha_r(t)=\sat(\alpha_{\mathrm{safe}}(t),\alpha_{\min},\alpha_{\max}),
  \label{eq:ch4_alpha_r}
\end{equation}
\begin{equation}
  \dot\alpha_{\mathrm{FP}}(t)=\lambda_{\alpha}(\alpha_r(t)-\alpha_{\mathrm{FP}}(t)),\quad \lambda_\alpha>0 .
  \label{eq:ch4_alpha_ode}
\end{equation}

\begin{lemma}[力--位优先级的区间不变性与变化率上界]
\label{lem:ch4_alpha_bound}
若 $\alpha_{\mathrm{FP}}(0)\in[\alpha_{\min},\alpha_{\max}]$，且 $\alpha_r(t)\in[\alpha_{\min},\alpha_{\max}]$，则\eqref{eq:ch4_alpha_ode}生成的 $\alpha_{\mathrm{FP}}(t)$ 满足
\begin{equation}
  \alpha_{\mathrm{FP}}(t)\in[\alpha_{\min},\alpha_{\max}],\quad \forall t\ge0,
  \label{eq:ch4_alpha_interval}
\end{equation}
并且
\begin{equation}
  |\dot\alpha_{\mathrm{FP}}(t)|\le \lambda_\alpha(\alpha_{\max}-\alpha_{\min}).
  \label{eq:ch4_alpha_rate}
\end{equation}
\end{lemma}

\begin{proof}
当 $\alpha_{\mathrm{FP}}=\alpha_{\min}$ 时，由 $\alpha_r(t)\ge\alpha_{\min}$ 可得 $\dot\alpha_{\mathrm{FP}}\ge0$；当 $\alpha_{\mathrm{FP}}=\alpha_{\max}$ 时，由 $\alpha_r(t)\le\alpha_{\max}$ 可得 $\dot\alpha_{\mathrm{FP}}\le0$。因此区间 $[\alpha_{\min},\alpha_{\max}]$ 为正不变集。进一步，由\eqref{eq:ch4_alpha_ode}有
\begin{equation}
  |\dot\alpha_{\mathrm{FP}}|=\lambda_\alpha|\alpha_r-\alpha_{\mathrm{FP}}|\le\lambda_\alpha(\alpha_{\max}-\alpha_{\min}).
\end{equation}
引理得证。
\end{proof}

\subsection{远心运动约束执行}
\label{subsec:ch4_rcm_mapping}

本章采用工具期望到法兰期望的正向映射。给定工具尖端期望位置 $x_t^d$，定义
\begin{equation}
  n_d=\frac{p_c-x_t^d}{\norm{p_c-x_t^d}},
  \label{eq:ch4_nd}
\end{equation}
法兰期望位置为
\begin{equation}
  x_F^d=x_t^d+Ln_d .
  \label{eq:ch4_xf}
\end{equation}
若 $r=\norm{p_c-x_t^d}$，则速度映射为
\begin{equation}
  \dot x_F^d=\left[I-\frac{L}{r}(I-n_dn_d^{\T})\right]\dot x_t^d .
  \label{eq:ch4_vf}
\end{equation}
其中 $I-n_dn_d^{\T}$ 为垂直于工具轴方向的投影矩阵。控制器在法兰空间计算位置误差和速度误差，在力维度使用工具接触力误差和力误差积分状态，随后生成法兰平移控制输入
\begin{equation}
  u_F=-K(\alpha_{\mathrm{FP}},\hat K_e)\xi_F,
  \label{eq:ch4_uF}
\end{equation}
其中 $\xi_F$ 为由法兰位置误差、法兰速度误差、工具接触力误差和力误差积分构成的增广状态。姿态控制输入 $u_\phi$ 与 $u_F$ 合成为广义力，关节力矩为
\begin{equation}
  \tau=J_F^{\T}\begin{bmatrix}u_F&u_\phi^{\T}\end{bmatrix}^{\T} .
  \label{eq:ch4_tau}
\end{equation}
该方式避免了先在工具空间产生控制力再进行杠杆缩放的物理歧义，也便于将 远心运动姿态调节纳入同一法兰空间闭环。
\section{鲁棒性与稳定性分析}
\label{sec:ch4_theory}

由于柔性接触任务不可避免地存在模型残差、刚度失配和在线仲裁变化，闭环分析需要同时覆盖名义最优性与扰动有界性。以下先给出名义折扣系统的稳定结论，再讨论残差项和优先级变化对稳定裕度的影响。

\subsection{名义稳定性与扰动有界性}
\label{subsec:ch4_nominal_theory}

\begin{theorem}[固定力--位优先级下的闭环稳定性]
\label{thm:ch4_fixed_stability}
在\cref{ass:ch4_stabilizable}成立的条件下，给定固定 $\alpha_{\mathrm{FP}}\in[\alpha_{\min},\alpha_{\max}]$，采用\cref{thm:ch4_discount_lqr}得到的反馈律后，若闭环矩阵
\begin{equation}
  A_{\mathrm{cl},\alpha}=A_c-B_cR_\alpha^{-1}B_c^{\T}P_\alpha
  \label{eq:ch4_Acl}
\end{equation}
为赫尔维茨，则标称闭环系统渐近稳定。
\end{theorem}

\begin{proof}
取李雅普诺夫函数 $V_\alpha(\xi)=\xi^{\T}P_\alpha\xi$。沿标称闭环系统 $\dot\xi=A_{\mathrm{cl},\alpha}\xi$ 求导，有
\begin{equation}
  \dot V_\alpha=\xi^{\T}(A_{\mathrm{cl},\alpha}^{\T}P_\alpha+P_\alpha A_{\mathrm{cl},\alpha})\xi .
\end{equation}
由折扣黎卡提方程\eqref{eq:ch4_discount_are}和反馈律定义可得
\begin{equation}
  A_{\mathrm{cl},\alpha}^{\T}P_\alpha+P_\alpha A_{\mathrm{cl},\alpha}= -Q_\alpha-P_\alpha B_cR_\alpha^{-1}B_c^{\T}P_\alpha+\gamma P_\alpha .
  \label{eq:ch4_discount_lyap_identity}
\end{equation}
由于折扣项 $\gamma P_\alpha$ 会减弱普通李雅普诺夫耗散不等式，实际闭环稳定性由稳定化黎卡提解和 $A_{\mathrm{cl},\alpha}$ 的赫尔维茨性保证。更具体地，若 $A_{\mathrm{cl},\alpha}$ 为赫尔维茨，则对任意 $W_\alpha\succ0$，存在唯一 $\Pi_\alpha\succ0$ 满足
\begin{equation}
  A_{\mathrm{cl},\alpha}^{\T}\Pi_\alpha+\Pi_\alpha A_{\mathrm{cl},\alpha}=-W_\alpha .
\end{equation}
取 $\bar V_\alpha=\xi^{\T}\Pi_\alpha\xi$，沿闭环系统求导得到
\begin{equation}
  \dot{\bar V}_\alpha=-\xi^{\T}W_\alpha\xi<0,
  \quad \xi\ne0 .
\end{equation}
因此标称闭环系统渐近稳定。定理得证。
\end{proof}

\paragraph{残差与参数失配分析。}
\label{subsec:ch4_residual_theory}

\begin{theorem}[残差项影响下的实际性能上界]
\label{thm:ch4_residual_performance}
考虑真实系统
\begin{equation}
  \dot\xi=A_c\xi+B_cu+d(\xi,t),
  \label{eq:ch4_real_system}
\end{equation}
其中残差项在局部紧集内满足 $\norm{d(\xi,t)}\le\bar d$。设 $u_\alpha^*(\xi)$ 为标称系统的折扣最优反馈律，并假设真实闭环轨迹满足
\begin{equation}
  \int_0^\infty \e^{-\gamma t}\norm{\xi(t)}^2\dd t\le c_\xi\norm{\xi_0}^2+c_d\bar d^2 .
  \label{eq:ch4_state_energy_bound}
\end{equation}
则真实系统采用标称反馈律时，其折扣代价满足
\begin{equation}
  J_\alpha^{\mathrm{real}}(\xi_0,u_\alpha^*)\le V_\alpha(\xi_0)+c_1\norm{\xi_0}^2+c_2\bar d^2,
  \label{eq:ch4_real_cost_bound}
\end{equation}
其中 $c_1,c_2>0$ 与 $P_\alpha$、$\gamma$、残差界和闭环衰减性质有关。
\end{theorem}

\begin{proof}
令 $V_\alpha(\xi)=\xi^{\T}P_\alpha\xi$。沿真实系统\eqref{eq:ch4_real_system}求导，有
\begin{equation}
  \dot V_\alpha=\nabla V_\alpha^{\T}(A_c\xi+B_cu_\alpha^*)+\nabla V_\alpha^{\T}d(\xi,t) .
  \label{eq:ch4_Vdot_real}
\end{equation}
根据折扣哈密顿--雅可比--贝尔曼方程，在标称最优控制律下有
\begin{equation}
  \nabla V_\alpha^{\T}(A_c\xi+B_cu_\alpha^*)+\xi^{\T}Q_\alpha\xi+(u_\alpha^*)^{\T}R_\alpha u_\alpha^*-\gamma V_\alpha=0 .
  \label{eq:ch4_hjb_on_policy}
\end{equation}
将\eqref{eq:ch4_Vdot_real}代入\eqref{eq:ch4_hjb_on_policy}，得到
\begin{equation}
  \xi^{\T}Q_\alpha\xi+(u_\alpha^*)^{\T}R_\alpha u_\alpha^*=\gamma V_\alpha-\dot V_\alpha+\nabla V_\alpha^{\T}d(\xi,t) .
  \label{eq:ch4_stage_real}
\end{equation}
两边乘以 $\e^{-\gamma t}$，并利用
\begin{equation}
  \frac{\dd}{\dd t}\left(\e^{-\gamma t}V_\alpha(\xi(t))\right)=\e^{-\gamma t}\dot V_\alpha-\gamma\e^{-\gamma t}V_\alpha,
\end{equation}
可得
\begin{align}
  &\e^{-\gamma t}\left[\xi^{\T}Q_\alpha\xi+(u_\alpha^*)^{\T}R_\alpha u_\alpha^*\right] \\
  &\quad =-\frac{\dd}{\dd t}\left(\e^{-\gamma t}V_\alpha(\xi(t))\right)+\e^{-\gamma t}\nabla V_\alpha^{\T}d(\xi,t).
\end{align}
对 $[0,T]$ 积分，得到
\begin{align}
  &\int_0^T\e^{-\gamma t}\left[\xi^{\T}Q_\alpha\xi+(u_\alpha^*)^{\T}R_\alpha u_\alpha^*\right]\dd t \\
  &\quad =V_\alpha(\xi_0)-\e^{-\gamma T}V_\alpha(\xi(T))+\int_0^T\e^{-\gamma t}\nabla V_\alpha^{\T}d(\xi,t)\dd t .
\end{align}
由于 $V_\alpha\ge0$，舍去非正项 $-\e^{-\gamma T}V_\alpha(\xi(T))$ 并令 $T$ 趋于无穷，得到
\begin{equation}
  J_\alpha^{\mathrm{real}}(\xi_0,u_\alpha^*)\le V_\alpha(\xi_0)+\int_0^\infty\e^{-\gamma t}|\nabla V_\alpha^{\T}d(\xi,t)|\dd t .
  \label{eq:ch4_real_cost_prebound}
\end{equation}
由于 $\nabla V_\alpha=2P_\alpha\xi$，有
\begin{equation}
  |\nabla V_\alpha^{\T}d(\xi,t)|\le2\norm{P_\alpha}\norm{\xi(t)}\bar d .
\end{equation}
对任意 $\varepsilon_d>0$，由杨氏不等式得到
\begin{equation}
  2\norm{P_\alpha}\norm{\xi(t)}\bar d\le \varepsilon_d\norm{\xi(t)}^2+\frac{\norm{P_\alpha}^2}{\varepsilon_d}\bar d^2 .
\end{equation}
代入\eqref{eq:ch4_real_cost_prebound}，并利用\eqref{eq:ch4_state_energy_bound}以及 $\int_0^\infty \e^{-\gamma t}\dd t=1/\gamma$，可得
\begin{equation}
  J_\alpha^{\mathrm{real}}(\xi_0,u_\alpha^*)\le V_\alpha(\xi_0)+\varepsilon_d c_\xi\norm{\xi_0}^2+\left(\varepsilon_d c_d+\frac{\norm{P_\alpha}^2}{\varepsilon_d\gamma}\right)\bar d^2 .
\end{equation}
令 $c_1=\varepsilon_d c_\xi$，$c_2=\varepsilon_d c_d+\norm{P_\alpha}^2/(\varepsilon_d\gamma)$，得到\eqref{eq:ch4_real_cost_bound}。定理得证。
\end{proof}

\begin{theorem}[接触模型参数误差下的最优反馈误差界]
\label{thm:ch4_param_mismatch}
设真实接触参数为 $\theta_{\mathrm{true}}=\theta_{\mathrm{nom}}+\Delta\theta$，其中
\begin{equation}
  \theta=\begin{bmatrix}\vecop(K_e)^{\T}&\vecop(B_e)^{\T}\end{bmatrix}^{\T} .
\end{equation}
对应系统矩阵满足
\begin{equation}
  A_{\mathrm{true}}=A_{\mathrm{nom}}+\Delta A,
  \quad \norm{\Delta A}\le L_A\norm{\Delta\theta} .
  \label{eq:ch4_A_mismatch}
\end{equation}
若标称折扣黎卡提方程存在稳定化解 $P_{\mathrm{nom}}$，且对应闭环矩阵具有稳定裕度，则当 $\norm{\Delta\theta}$ 足够小时，真实黎卡提解 $P_{\mathrm{true}}$ 存在，并满足
\begin{equation}
  \norm{P_{\mathrm{true}}-P_{\mathrm{nom}}}\le C_P\norm{\Delta\theta} .
  \label{eq:ch4_P_error}
\end{equation}
进一步地，最优反馈增益和最优控制输入满足
\begin{equation}
  \norm{K_{\mathrm{true}}-K_{\mathrm{nom}}}\le C_K\norm{\Delta\theta},
  \label{eq:ch4_K_error}
\end{equation}
\begin{equation}
  \norm{u_{\mathrm{true}}^*(\xi)-u_{\mathrm{nom}}^*(\xi)}\le C_K\norm{\Delta\theta}\norm{\xi} .
  \label{eq:ch4_u_error}
\end{equation}
\end{theorem}

\begin{proof}
定义折扣黎卡提算子
\begin{equation}
  \mathcal{R}(P,A)=A^{\T}P+PA-PB_cR_\alpha^{-1}B_c^{\T}P+Q_\alpha-\gamma P .
\end{equation}
标称解满足 $\mathcal{R}(P_{\mathrm{nom}},A_{\mathrm{nom}})=0$，真实解满足 $\mathcal{R}(P_{\mathrm{true}},A_{\mathrm{true}})=0$。令
\begin{equation}
  \Delta P=P_{\mathrm{true}}-P_{\mathrm{nom}},
  \quad \Delta A=A_{\mathrm{true}}-A_{\mathrm{nom}} .
\end{equation}
在 $(P_{\mathrm{nom}},A_{\mathrm{nom}})$ 处展开，得到
\begin{equation}
  0=\mathcal{L}(\Delta P)+\Delta A^{\T}P_{\mathrm{nom}}+P_{\mathrm{nom}}\Delta A+\mathcal{N}(\Delta P,\Delta A),
  \label{eq:ch4_riccati_perturb}
\end{equation}
其中线性算子为
\begin{equation}
  \mathcal{L}(\Delta P)=A_{\mathrm{cl,nom}}^{\T}\Delta P+\Delta P A_{\mathrm{cl,nom}}-\gamma\Delta P,
\end{equation}
而 $\mathcal{N}(\Delta P,\Delta A)$ 包含二阶小量，例如 $\Delta A^{\T}\Delta P$、$\Delta P\Delta A$ 和 $-\Delta P B_cR_\alpha^{-1}B_c^{\T}\Delta P$。由于 $A_{\mathrm{cl,nom}}$ 具有稳定裕度，$\mathcal{L}$ 可逆，并存在常数 $c_L>0$，使得
\begin{equation}
  \norm{\mathcal{L}^{-1}(Y)}\le c_L\norm{Y} .
\end{equation}
由\eqref{eq:ch4_riccati_perturb}可得
\begin{equation}
  \norm{\Delta P}\le c_L\left(2\norm{P_{\mathrm{nom}}}\norm{\Delta A}+\norm{\mathcal{N}(\Delta P,\Delta A)}\right).
  \label{eq:ch4_deltaP_pre}
\end{equation}
当 $\norm{\Delta A}$ 和 $\norm{\Delta P}$ 足够小时，存在常数 $c_N>0$，使得
\begin{equation}
  \norm{\mathcal{N}(\Delta P,\Delta A)}\le c_N\left(\norm{\Delta A}\norm{\Delta P}+\norm{\Delta P}^2\right).
\end{equation}
利用小扰动条件，可将二阶项吸收到\eqref{eq:ch4_deltaP_pre}左端，得到
\begin{equation}
  \norm{\Delta P}\le C_A\norm{\Delta A} .
\end{equation}
再由\eqref{eq:ch4_A_mismatch}得到\eqref{eq:ch4_P_error}，其中 $C_P=C_A L_A$。由于
\begin{equation}
  K_{\mathrm{true}}-K_{\mathrm{nom}}=R_\alpha^{-1}B_c^{\T}(P_{\mathrm{true}}-P_{\mathrm{nom}}),
\end{equation}
有
\begin{equation}
  \norm{K_{\mathrm{true}}-K_{\mathrm{nom}}}\le \norm{R_\alpha^{-1}B_c^{\T}}\,C_P\norm{\Delta\theta} .
\end{equation}
令 $C_K=\norm{R_\alpha^{-1}B_c^{\T}}C_P$，得到\eqref{eq:ch4_K_error}。最后，由 $u^*(\xi)=-K\xi$ 可得\eqref{eq:ch4_u_error}。定理得证。
\end{proof}

\paragraph{仲裁速率与稳定裕度。}
\label{subsec:ch4_rate_margin}

动态 $\alpha_{\mathrm{FP}}(t)$ 提高了控制器对力风险的响应能力，但也会引入时变反馈增益。为刻画其影响，定义在线近似值函数
\begin{equation}
  \hat V_{\alpha}(\xi)=\xi^{\T}P_{\alpha}\xi .
  \label{eq:ch4_val_online}
\end{equation}
若在线 $P_{\alpha}$ 由离线表插值得到，或者由少量网格点组合而成，则可能存在贝尔曼残差
\begin{equation}
  \Delta_\alpha=A_c^{\T}P_\alpha+P_\alpha A_c-P_\alpha B_cR_\alpha^{-1}B_c^{\T}P_\alpha+Q_\alpha-\gamma P_\alpha .
  \label{eq:ch4_bellman_residual}
\end{equation}
令
\begin{equation}
  \epsilon_B=\max_{\alpha\in[\alpha_{\min},\alpha_{\max}]}\norm{\Delta_\alpha} .
  \label{eq:ch4_eps_B}
\end{equation}
当 $P_\alpha$ 是精确折扣 ARE 解时，$\epsilon_B=0$；当存在插值误差、模型近似或数值误差时，$\epsilon_B$ 表示值函数近似误差的上界。

\begin{theorem}[时变优先级下的性能损失上界]
\label{thm:ch4_timevarying_performance}
设在线策略 $\hat u_\alpha(\xi)=-R_\alpha^{-1}B_c^{\T}P_\alpha\xi$ 可容许，且存在常数 $c>0$，使得
\begin{equation}
  \int_0^\infty \e^{-\gamma t}\norm{\xi(t)}^2\dd t\le c\norm{\xi_0}^2 .
\end{equation}
若 $\norm{\Delta_\alpha}\le\epsilon_B$，$\norm{\partial P_\alpha/\partial\alpha}\le L_P$，且 $|\dot\alpha_{\mathrm{FP}}(t)|\le v_\alpha$，则相对于理想时变折扣最优策略，在线近似策略的性能损失满足
\begin{equation}
  J_{\alpha(\cdot)}(\xi_0;\hat u_{\alpha(\cdot)})-J_{\alpha(\cdot)}^*(\xi_0)
  \le 2c(\epsilon_B+v_\alpha L_P)\norm{\xi_0}^2 .
  \label{eq:ch4_perf_loss}
\end{equation}
\end{theorem}

\begin{proof}
对时变近似值函数 $\hat V_{\alpha(t)}(\xi)=\xi^{\T}P_{\alpha(t)}\xi$，有
\begin{equation}
  \frac{\dd}{\dd t}\left(\e^{-\gamma t}\hat V_{\alpha(t)}(\xi(t))\right)=\e^{-\gamma t}\left[\nabla \hat V_\alpha^{\T}\dot\xi-\gamma\hat V_\alpha+\frac{\partial \hat V_\alpha}{\partial\alpha}\dot\alpha\right].
\end{equation}
由哈密顿函数定义可得，对任意可容许控制 $u$，
\begin{equation}
  J_{\alpha(\cdot)}(\xi_0;u)=\hat V_{\alpha(0)}(\xi_0)+\int_0^\infty\e^{-\gamma t}\left[H_\alpha(\xi,\nabla\hat V_\alpha,u)+\frac{\partial \hat V_\alpha}{\partial\alpha}\dot\alpha\right]\dd t,
  \label{eq:ch4_J_identity_timevarying}
\end{equation}
其中终端项在稳定性和折扣条件下消失。对在线贪婪策略 $\hat u_\alpha$，哈密顿函数残差满足
\begin{equation}
  |H_\alpha(\xi,\nabla\hat V_\alpha,
  \hat u_\alpha)|=|\xi^{\T}\Delta_\alpha\xi|\le\epsilon_B\norm{\xi}^2 .
\end{equation}
同时
\begin{equation}
  \left|\frac{\partial \hat V_\alpha}{\partial\alpha}\dot\alpha\right|\le L_Pv_\alpha\norm{\xi}^2 .
\end{equation}
因此在线策略的代价上界为
\begin{equation}
  J_{\alpha(\cdot)}(\xi_0;\hat u_{\alpha(\cdot)})\le \hat V_{\alpha(0)}(\xi_0)+c(\epsilon_B+v_\alpha L_P)\norm{\xi_0}^2 .
\end{equation}
另一方面，理想最优策略的哈密顿函数不小于在线贪婪策略对应的最小哈密顿函数，因此可得下界
\begin{equation}
  J_{\alpha(\cdot)}^*(\xi_0)\ge \hat V_{\alpha(0)}(\xi_0)-c(\epsilon_B+v_\alpha L_P)\norm{\xi_0}^2 .
\end{equation}
将上下界相减，即得到\eqref{eq:ch4_perf_loss}。定理得证。
\end{proof}

\begin{remark}
\cref{thm:ch4_timevarying_performance}说明，优先级变化率 $v_\alpha$ 不是单纯的工程滤波参数，而是影响性能损失上界的设计变量。较大的 $v_\alpha$ 提高力风险响应速度，但会增大 $v_\alpha L_P$ 项；较小的 $v_\alpha$ 能减小性能损失界，但可能使力越界响应滞后。因此，$\lambda_\alpha$ 应根据任务安全裕度、采样频率和闭环稳定裕度共同选择。
\end{remark}

\section{仿真结果与分析}
\label{sec:ch4_simulation}

\subsection{仿真设置}
\label{subsec:ch4_sim_task}

仿真任务设置为柔性表面恒力扫描。工具尖端沿给定方向低速移动，并与分段变刚度环境保持连续接触。为突出第四章方法相对于第三章刚性几何约束共享控制的增量，本章分别设置无远心运动约束和有远心运动约束两类场景。前者直接考察力--位仲裁策略在接触扫描中的作用；后者进一步引入长工具几何映射和远心运动约束误差，考察控制器在接触力调节、扫描轨迹跟踪和通道合规之间的综合权衡。典型参数如\cref{tab:ch4_env_params}所示。该设置用于检验控制器在刚度突变时是否能够降低力超调，并在风险降低后恢复位置跟踪性能。

\begin{table}[H]
\centering
\caption{第四章柔性接触扫描仿真环境参数}
\label{tab:ch4_env_params}
\begin{tabular}{llll}
\toprule
场景 & 扫描区间与期望力 & 刚度 $K_e$ & 阻尼 $B_e$ \\
\midrule
无远心运动 & $x:0.40\to0.48\,\mathrm{m}$, $F_d=1.0\,\mathrm{N}$ & $300/500\,\mathrm{N/m}$ & $5/8\,\mathrm{N\,s/m}$ \\
远心运动 & $x:0.23\to0.32\,\mathrm{m}$, $F_d=0.5\,\mathrm{N}$ & $250/500\,\mathrm{N/m}$ & $4/8\,\mathrm{N\,s/m}$ \\
远心运动几何 & $p_c=[0.3,0,0.235]^{\T}\,\mathrm{m}$, $L=0.525\,\mathrm{m}$ & -- & -- \\
\bottomrule
\end{tabular}
\end{table}

对比方法包括三组：无远心运动在线优先级仲裁、无远心运动连续力安全裕度仲裁，以及远心运动连续力安全裕度仲裁。前两组用于观察在不施加几何约束时力--位仲裁本身的作用；第三组进一步加入远心运动几何映射和约束误差记录，用于检验第四章方法相对于第三章刚性约束共享控制的增量，即在保持工具通道合规的同时处理柔性接触力安全问题。

\paragraph{评价指标与参数设置。}
\label{subsec:ch4_metrics}

力跟踪误差指标为
\begin{equation}
  E_F^{\mathrm{RMS}}=\sqrt{\frac{1}{N}\sum_{k=1}^{N}(y_{f,k}-F_d)^2} .
  \label{eq:ch4_force_rmse}
\end{equation}
力安全边界违反次数为
\begin{equation}
  N_{\mathrm{upper}}=\sum_{k=1}^{N}\mathbb{I}(y_{f,k}>F_{\max}),
  \quad
  N_{\mathrm{lower}}=\sum_{k=1}^{N}\mathbb{I}(y_{f,k}<F_{\min}) .
  \label{eq:ch4_violation_count}
\end{equation}
位置跟踪误差为
\begin{equation}
  E_x^{\mathrm{RMS}}=\sqrt{\frac{1}{N}\sum_{k=1}^{N}\norm{x_k-x_{d,k}}^2} .
  \label{eq:ch4_position_rmse}
\end{equation}
几何约束误差为
\begin{equation}
  E_{\mathrm{gc}}^{\mathrm{RMS}}=\sqrt{\frac{1}{N}\sum_{k=1}^{N}e_{\mathrm{gc},k}^2} .
  \label{eq:ch4_gc_rmse}
\end{equation}
控制输入强度用平均关节力矩范数表示：
\begin{equation}
  \bar\tau=\frac{1}{N}\sum_{k=1}^{N}\norm{\tau_k} .
  \label{eq:ch4_tau_metric}
\end{equation}
此外，还应记录单步计算时间和 $\alpha_{\mathrm{FP}}(t)$ 变化曲线，以验证离线--在线结构的实时性和动态仲裁的可解释性。

\begin{table}[H]
\centering
\caption{第四章柔性接触扫描仿真实验结果}
\label{tab:ch4_result_template}
\scriptsize
\begin{tabular}{lccc}
\toprule
指标 & 无远心运动在线优先级 & 无远心运动连续力裕度 & 远心运动连续力裕度 \\
\midrule
采样数/持续时间 & $4000/40.09\,\mathrm{s}$ & $4000/40.13\,\mathrm{s}$ & $9628/96.57\,\mathrm{s}$ \\
力均方根误差/N & $0.126$ & $0.114$ & $0.095$ \\
力平均绝对误差/N & $0.073$ & $0.065$ & $0.052$ \\
力峰值绝对误差/N & $0.990$ & $0.960$ & $0.968$ \\
上界/下界违反次数 & $0/18$ & $0/3$ & $62/152$ \\
位置均方根误差/mm & $4.00$ & $4.40$ & $4.90$ \\
位置峰值误差/mm & $6.50$ & $10.41$ & $7.20$ \\
几何约束均方根误差/mm & $0$ & $0$ & $4.85$ \\
几何约束峰值误差/mm & $0$ & $0$ & $7.20$ \\
仲裁权重均值/范围 & $0.789/(0.170,0.944)$ & $0.641/(0.250,0.910)$ & $0.700/(0.250,0.872)$ \\
控制输入均方根 & $0.551$ & $0.929$ & $0.300$ \\
\bottomrule
\end{tabular}
\end{table}

\begin{figure}[H]
\centering
\includegraphics[width=\textwidth]{figures/ch3_ch4/ch4_force_position_rcm_compare.png}
\caption{第四章柔性接触扫描中力误差、位置误差、仲裁权重和远心运动约束误差的时序比较}
\label{fig:ch4_force_position_rcm_compare}
\end{figure}

\subsection{结果分析}
\label{subsec:ch4_results_analysis}

柔性接触仿真从三条主线展开分析。第一是力安全。无远心运动场景中，连续力安全裕度仲裁相对于在线优先级仲裁将力均方根误差由 $0.126\,\mathrm{N}$ 降至 $0.114\,\mathrm{N}$，下界违反次数由 $18$ 次降至 $3$ 次，说明基于力裕度的连续仲裁能够更直接地抑制接触力偏离。其位置均方根误差从 $4.00\,\mathrm{mm}$ 增至 $4.40\,\mathrm{mm}$，体现了力安全优先时对位置跟踪精度的主动让渡。\cref{fig:ch4_force_position_rcm_compare}中仲裁权重的变化也显示，当接触力出现突变或进入高刚度区域时，动态仲裁会降低位置优先级并增强力调节作用，随后在力误差减小时逐渐恢复位置跟踪。

第二是远心运动合规。加入远心运动几何映射后，控制器需要同时维持接触力、扫描轨迹和通道约束；最终运行的几何约束均方根误差为 $4.85\,\mathrm{mm}$，位置均方根误差为 $4.90\,\mathrm{mm}$，力均方根误差为 $0.095\,\mathrm{N}$。该结果表明，在长工具约束带来的额外自由度耦合下，连续力裕度仲裁仍能保持较低的力误差，但需要付出有限的几何约束误差代价。与无远心运动场景相比，有远心运动场景的实验时间更长、约束项更多，因此不能仅以单一误差最小作为评价标准，而应关注力误差、位置误差和几何约束误差三者是否维持在可接受范围内。

第三是与第三章仿真的互补关系。第三章结果说明，合作博弈控制和模糊自适应仲裁能够降低刚性几何约束任务中的远心运动误差和末端跟踪误差；但第三章实验的主要风险来自人机意图切换和障碍接近，并未把持续接触力作为状态反馈目标。第四章在此基础上加入接触力误差、泄漏积分、环境刚度估计和力安全边界，使仲裁变量从人机角色分配扩展到力--位任务分配。因而，第三章提供了几何通道约束与动态仲裁的基线证据，第四章则补充了柔性接触和模型失配条件下的力安全验证。

\section{本章小结}
\label{sec:ch4_summary}

本章针对柔性接触恒力扫描任务，提出了一种交互力安全裕度驱动的折扣帕累托力--位协同控制方法。首先，从修正阻抗动力学和开尔文--沃伊特局部接触模型出发，构造了包含位置误差、速度误差、力误差和力误差泄漏积分的增广状态，并将真实系统表示为含残差项的非线性增广误差系统。随后，本章结合表观刚度低通估计和有界递推最小二乘阻尼估计，给出了用于增益调度的环境参数在线估计方法。在此基础上，本章先在折扣型位置目标和力目标代价函数下建模并求解帕累托博弈控制器，再描述交互力安全裕度驱动的动态仲裁方案。离线求解、在线查表和显式反馈的结构满足实时接触控制需求。

在仲裁机制方面，本章将接触力相对于上下安全边界的归一化裕度作为核心风险指标，设计了动态力--位优先级 $\alpha_{\mathrm{FP}}(t)$。该参数不同于第三章的人机仲裁参数，其作用是在位置跟踪和力调节之间进行任务级权衡。理论分析表明，在标称模型下，折扣帕累托控制律具有明确的最优性；在存在残差扰动时，真实性能相对标称性能的偏差可由残差界刻画；在环境刚度和阻尼建模不准确时，黎卡提解、反馈增益和最优控制输入相对真实最优解的误差均具有有界性。

至此，第三章和第四章形成递进关系。第三章解决刚性约束任务中的人机意图仲裁与共享控制问题，第四章进一步处理柔性连续接触中的力--位协调和模型不确定性问题。二者共同构成面向约束操作和接触交互任务的分层合作博弈控制方法基础。
